\sect{मार्गशीर्ष-कृष्ण-उत्पन्ना-एकादशी-कृत-मुरवधः}
\label{sec:padma-margashirsha-krishna-utpanna}

\ganapatyadiDhyanam
\padmaPuranam

\dnsub{अथ चत्वारिंशोऽध्यायः॥४०}
\uvacha{महादेव उवाच}

\twolineshloka
{एकस्मिन् समये पुत्र गतोऽहं विष्णुसन्निधौ}
{तत्र पृष्टं मया पूर्वं माहात्म्यं द्वादशीभवम्}% ॥१॥

\uvacha{नारद उवाच}

\twolineshloka
{कीदृशी स्यान्महादेव महाद्वादशिका परा}
{तस्या व्रते फलं कीदृग्वद सर्वेश्वर प्रभो}% ॥२॥

\uvacha{शिव उवाच}

\twolineshloka
{इयमेकादशी ब्रह्मन् महापुण्यफलप्रदा}
{ऋक्षयोगैश्च संयुक्ता कर्तव्या मुनिसत्तमैः}% ॥३॥

\twolineshloka
{जया च विजया चैव जयन्ती पापनाशिनी}
{सर्वपापहराश्चैताः कर्तव्याः फलकाङ्क्षिभिः}% ॥४॥

\twolineshloka
{एकादश्यां यदा ऋक्षं शुक्लपक्षे पुनर्वसुः}
{नाम्ना सा च जया ख्याता तिथीनामुत्तमा तिथिः}% ॥५॥

\twolineshloka
{तामुपोष्य नरः पापान्मुच्यते नात्र संशयः}
{यदा च शुक्लद्वादश्यां नक्षत्रं श्रवणं भवेत्}% ॥६॥

\twolineshloka
{विजया सा समाख्याता तिथीनामुत्तमा तिथिः}
{सहस्रगुणितं दानं यस्यां वै विप्र भोजनम्}% ॥७॥

\twolineshloka
{होमस्तथोपवासश्च सहस्राधिफलप्रदः}
{यदा च शुक्लद्वादश्यां प्राजापत्यं हि जायते}% ॥८॥

\twolineshloka
{जयन्ती नाम सा प्रोक्ता सर्वपापहरा तिथिः}
{सप्तजन्मकृतं पापं स्वल्पं वा यदि वा बहु}% ॥९॥

\twolineshloka
{प्रक्षालयति गोविन्दस्तस्यामभ्यर्चितो ध्रुवम्}
{यदा वै शुक्लद्वादश्यां पुष्यं भवति कर्हिचित्}% ॥१०॥

\twolineshloka
{तदा तु सा महापुण्या भविता पापनाशिनी}
{यो ददाति तिलप्रस्थं नित्यं संवत्सरं प्रति}% ॥११॥

\twolineshloka
{उपवासं च यस्तस्यां करोत्येतत् समं स्मृतम्}
{तस्यां जगत्पतिर्देवस्तुष्टः सर्वेश्वरो हरिः}% ॥१२॥

\twolineshloka
{प्रत्यक्षतां प्रयात्येव तत्रानन्तफलं स्मृतम्}
{सगरेण ककुत्स्थेन नहुषेण च साधितः}% ॥१३॥

\twolineshloka
{तस्यामाराधितः कृष्णो दत्तवानखिलं भुवि}
{वाचिकान्मानसाद्वाऽपि कायजाच्च विशेषतः}% ॥१४॥

\twolineshloka
{सप्तजन्मकृतात् पापान्मुच्यते नात्र संशयः}
{तामेकां समुपोष्याथ पुण्यनक्षत्रसंयुताम्}% ॥१५॥

\twolineshloka
{एकादशीसहस्रस्य फलं प्राप्नोति मानवः}
{स्नानं दानं जपो होमः स्वाध्यायो देवतार्चनम्}% ॥१६॥

\twolineshloka
{यत् तस्यां क्रियते किञ्चित्तदक्षयफलं स्मृतम्}
{तस्मादेषा प्रकर्तव्या यत्नेन फलकाङ्क्षिभिः}% ॥१७॥

\twolineshloka
{पञ्चमेनाश्वमेधेन यदि स्नातो युधिष्ठिरः}
{पर्यपृच्छत धर्मात्मा कृष्णं यदुकुलोद्वहम्}% ॥१८॥

\uvacha{युधिष्ठिर उवाच}

\twolineshloka
{उपवासस्य नक्तस्य त्वेकभुक्तस्य मे प्रभो}
{किं पुण्यं किं फलं तस्य ब्रूहि सर्वं जनार्दन}% ॥१९॥

\uvacha{श्रीभगवानुवाच}
\onelineshloka*
{हेमन्ते चैव सम्प्राप्ते मासे मार्गे च शोभने}

\twolineshloka
{कृष्णपक्षे च या पार्थ द्वादशी तामुपोषयेत्}
{दशम्यां चैकभक्तश्च शुद्धचित्तो दृढव्रतः}% ॥२०॥

\twolineshloka
{नक्तं चैव तथा ज्ञात्वा दशम्यां नियतः सदा}
{दिवसस्याष्टमे भागे मन्दीभूते दिवाकरे}% ॥२१॥

\twolineshloka
{नक्तं च तद् विजानीयान्न नक्तं निशिभोजनम्}
{नक्षत्रदर्शनान्नक्तं गृहस्थस्य विधीयते}% ॥२२॥

\twolineshloka
{यतेर्दिनाष्टमे भागे रात्रौ तस्य निषेधनम्}
{ततः प्रभातसमये कृत्वा च नियमं व्रती}% ॥२३॥

\twolineshloka
{मध्याह्ने च तथा पार्थ स्नानं शुचिः समाचरेत्}
{अधमं कूपके स्नानं वाप्यां स्नानं च मध्यमम्}% ॥२४॥

\twolineshloka
{तडागे चोत्तमं स्नानं नद्यां स्नानं ततः परम्}
{पीड्यन्ते जन्तवो यत्र जलमध्ये व्यवस्थिते}% ॥२५॥

\twolineshloka
{तत्र स्नाने कृते पार्थ पापं पुण्यं समं भवेत्}
{गृहे चैवोत्तमं स्नानं जलं चैव विशोधयेत्}% ॥२६॥

\twolineshloka
{तस्मात् तु पाण्डवश्रेष्ठ गृहे स्नानं समाचरेत्}
{अश्वक्रान्ते रथक्रान्ते विष्णुक्रान्ते वसुन्धरे}% ॥२७॥

\twolineshloka
{मृत्तिके हर मे पापं यन्मया पूर्वसञ्चितम्}
{क्रोधलोभौ परित्यज्य चैकचित्तो दृढव्रतः}% ॥२८॥

\twolineshloka
{नालापेतान्त्यजं चैव तथा पाखण्डिनो नरान्}
{मिथ्यावादरतांश्चैव तथा ब्राह्मणनिन्दकान्}% ॥२९॥

\twolineshloka
{अन्यांश्चैव दुराचारानगम्यागमनेरतान्}
{परद्रव्यापहारांश्च परदाराभिगामिनः}% ॥३०॥

\twolineshloka
{केशवं पूजयित्वा तु नैवेद्यं तत्र कारयेत्}
{दीपं दद्याद्गृहे तत्र भक्तियुक्तेन चेतसा}% ॥३१॥

\twolineshloka
{तद्दिने वर्जयेत् पार्थ निद्रां चैव तु मैथुनम्}
{धर्मशास्त्रविनोदेन दिनं सर्वं निवारयेत्}% ॥३२॥

\twolineshloka
{रात्रौ जागरणं कृत्वा भक्तियुक्तो नृपोत्तम}
{विप्रेभ्यो दक्षिणां दद्यात् प्रणिपत्य क्षमापयेत्}% ॥३३॥

\twolineshloka
{यथा कृष्णा तथा शुक्ला विधिनैवं प्रकारयेत्}
{एकादशीं द्विजः पार्थ विभेदं नैव कारयेत्}% ॥३४॥

\twolineshloka
{एवं हि कुरुते यस्तु शृणु तस्य हि यत्फलम्}
{शङ्खोद्धारे नरः स्नात्वा दृष्ट्वा देवं गदाधरम्}% ॥३५॥

\twolineshloka
{एकादश्युपवासस्य कलां नार्हति षोडशीम्}
{सङ्क्रान्तिषु चतुर्लक्षं यो ददाति नृपोत्तम}% ॥३६॥

\twolineshloka
{एकादश्युपवासस्य कलां नार्हति षोडशीम्}
{प्रभासक्षेत्रे यत् पुण्यं ग्रहणे चन्द्र सूर्ययोः}% ॥३७॥

\twolineshloka
{तत्फलं जायते नूनमेकादश्युपवासिनः}
{केदारे चोदकं पीत्वा पुनर्जन्म न विद्यते}% ॥३८॥

\twolineshloka
{तथा चैकादशी पार्थ गर्भवासक्षयङ्करी}
{अश्वमेधस्य यज्ञस्य पृथिव्यां यत्फलं लभेत्}% ॥३९॥

\twolineshloka
{तस्माच्छतगुणं पुण्यमेकादश्युपवासिनः}
{तपस्विनो गृहे यस्य भुञ्जते च द्विजोत्तमाः}% ॥४०॥

\twolineshloka
{तत्फलं जायते नूनमेकादश्युपवासिनः}
{गोसहस्रेण यत् पुण्यं दत्वा वेदान्तपारगे}% ॥४१॥

\twolineshloka
{तस्माच्छतगुणं पुण्यमेकादश्युपवासिनः}
{येषां देहे त्रयो देवा ब्रह्मविष्णुमहेश्वराः}% ॥४२॥

\twolineshloka
{वसन्ति तेषां ते तुल्या एकादश्युपवासिनः}
{ते नराः पुण्यकर्माणो ये भक्ता हरिपूजकाः}% ॥४३॥

\twolineshloka
{एकादशीव्रतस्यापि पुण्यसङ्ख्या न विद्यते}
{एतत् पुण्यं भवेत् तस्य यत्सुरैरपि दुर्लभम्}% ॥४४॥

\twolineshloka
{एतस्मादर्धपुण्यं तु प्राप्यते नक्तभोजनात्}
{नक्तस्यार्धं भवेत् पुण्यमेकभक्तेन देहिनाम्}% ॥४५॥

\twolineshloka
{तावद्गर्जन्ति तीर्थानि दानानि नियमानि च}
{यावन्नोपोषयेज्जन्तुर्वासरं विष्णुवल्लभम्}% ॥४६॥

\twolineshloka
{तस्मात् त्वं पाण्डवश्रेष्ठ व्रतमेतत् समाचर}
{पुण्यसङ्ख्यां न जानामि यत्त्वं पृच्छसि पाण्डव}% ॥४७॥

\twolineshloka
{एतद्धि कथितं पार्थ यद् गोप्यं व्रतमुत्तमम्}
{एकादशीसमं नास्ति कृत्वा यज्ञसहस्रकम्}% ॥४८॥

\uvacha{युधिष्ठिर उवाच}

\twolineshloka
{उत्पन्ना सा कथं देव पुण्या चैकादशी तिथिः}
{कथं पवित्रा विश्वेऽस्मिन् कथं वै देवताप्रिया}% ॥४९॥

\uvacha{श्रीभगवानुवाच}

\twolineshloka
{पुरा कृतयुगे पार्थ मुरनामेति दानवः}
{अत्यद्भुतो महारौद्र सर्वदेवभयङ्करः}% ॥५०॥

\twolineshloka
{इन्द्रोऽपि निर्जितस्तेन सर्वदेवास्तथा नृप}
{महासुरेण तेनैव मृत्युना च दुरात्मना}% ॥५१॥

\twolineshloka
{स्वर्गान्निराकृतास्तेन विचरन्ति महीतले}
{सशङ्का भयभीताश्च सर्वे गत्वा महेश्वरम्}% ॥५२॥

\twolineshloka
{इन्द्रेण कथितं सर्वमीश्वरस्यापि चाग्रतः}
{स्वर्गलोकपरिभ्रष्टा विचरन्ति  महीतले}% ॥५३॥

\twolineshloka
{मर्त्येषु संस्थिता देवा न शोभन्ते महेश्वर}
{उपायं ब्रूहि मे देव ह्यमरा यान्ति कां गतिम्}% ॥५४॥

\uvacha{महादेव उवाच}

\twolineshloka
{देवराज सुरश्रेष्ठ यत्राऽऽस्ते गरुडध्वजः}
{शरण्यश्च जगन्नाथः परित्राता परायणः}% ॥५५॥

\twolineshloka
{तत्र गच्छ सुरश्रेष्ठ स वो रक्षां विधास्यति}
{ईश्वरस्य वचः श्रुत्वा देवराजो महामतिः}% ॥५६॥

\twolineshloka
{त्रिदशैः सहितो यत्र गतस्तत्र युधिष्ठिर}
{जलमध्ये प्रसुप्तं तं दृष्ट्वा देवं गदाधरम्}% ॥५७॥

\onelineshloka
{कृताञ्जलिपुटो भूत्वा इन्द्रः स्तोत्रमुदीरयत्}% ॥५८॥

\uvacha{इन्द्र उवाच}

\twolineshloka
{ॐ नमो देवदेवेश देवदानववन्दित}
{दैत्यारे पुण्डरीकाक्ष त्राहि नो मधुसूदन}% ॥५९॥

\twolineshloka
{सुराः सर्वे समायाता भयभीताश्च दानवात्}
{शरणं त्वां जगन्नाथ त्राहि मां भक्तवत्सल}% ॥६०॥

\twolineshloka
{त्राहि नो देवदेवेश त्राहि त्राहि जनार्दन}
{त्राहि वै पुण्डरीकाक्ष दानवानां विनाशक}% ॥६१॥

\twolineshloka
{त्वत्समीपं गताः सर्वे त्वमेव शरणं प्रभो}
{शरणागतानां देवानां सहायं कुरु वै प्रभो}% ॥६२॥

\twolineshloka
{त्वं पतिस्त्वं मतिर्देव त्वं कर्ता त्वं च कारणम्}
{त्वं माता सर्वलोकानां त्वमेव जगतः पिता}% ॥६३॥

\twolineshloka
{भगवन् देवदेवेश शरणागतवत्सल}
{शरणं तव चाऽऽयाता भयभीताश्च देवताः}% ॥६४॥


\threelineshloka
{देवता निर्जिताः सर्वाः स्वर्गभ्रष्टाः कृताः प्रभो}
{अत्युग्रेण हि दैत्येन मुरनाम्ना महौजसा}
{इन्द्रस्य वचनं श्रुत्वा विष्णुर्वचनमब्रवीत्}% ॥६५॥

\uvacha{श्रीभगवानुवाच}

\threelineshloka
{कीदृशो दानवः शक्र किं रूपं कीदृशं बलम्}
{क्व स्थानं तस्य दुष्टस्य किं वीर्यं कः पराक्रमः}
{किं वरं तस्य दुष्टस्य ममाऽऽख्याहि महामते}% ॥६६॥

\uvacha{इन्द्र उवाच}

\twolineshloka
{पूर्वं बभूव देवेश ब्रह्मवंशसमुद्भवः}
{तालजङ्घस्तु नाम्ना च अत्युग्रोऽपि महासुरः}% ॥६७॥

\twolineshloka
{तस्य पुत्रो हि विख्यातो मुरनामेति दानवः}
{अत्युत्कटो महावीर्यो देवतानां भयङ्करः}% ॥६८॥

\twolineshloka
{पुरी चन्द्रावती नाम्ना तत्र स्थाने वसत्यसौ}
{निर्जिता देवताः सर्वा स्वर्गात्तेन विवासिताः}% ॥६९॥

\twolineshloka
{इन्द्रोऽन्यो रोपितस्तेन वातश्चैव हुताशनः}
{चन्द्रसूर्यो कृतौ चान्यौ वायुर्वरुण एव च}% ॥७०॥

\twolineshloka
{सर्वमात्मकृतं तेन सत्यं सत्यं जनार्दन}
{देवलोकः कृतस्तेन सर्वस्थानविवर्जितः}% ॥७१॥

\twolineshloka
{तस्य तद् वचनं श्रुत्वा कोपमानो जनार्दनः}
{हनिष्ये दानवं दुष्टं देवतानां भयङ्करम्}% ॥७२॥

\twolineshloka
{त्रिदशैः सहितो देवो गतश्चन्द्रावतीं पुरीम्}
{दृष्टो देवैस्तु दैत्येन्द्रो गर्जमानः पुनः पुनः}% ॥७३॥

\twolineshloka
{तेन सर्वे जिता देव गताश्चैव दिशो दश}
{हरिं निरीक्ष्य प्रोवाच तिष्ठ तिष्ठेति दानवः}% ॥७४॥

\onelineshloka*
{भगवानब्रवीत् तं च क्रोधसंरक्तलोचनः}

\uvacha{श्रीभगवानुवाच}

\onelineshloka
{रे दानव दुराचार मम बाहुं निरीक्षय}% ॥७५॥

\twolineshloka
{ततस्ते सम्मुखाः सर्वे विष्णुना दुष्टदानवाः}
{हता बाणैः पुनर्दिव्यैर्जाताश्च भयविह्वलाः}% ॥७६॥

\twolineshloka
{चक्रं मुक्तं च कृष्णेन दैत्यसैन्येषु पाण्डव}
{तेन छिन्नास्तु शतशो बहवो निधनं गताः}% ॥७७॥

\twolineshloka
{एकोऽपि दानवस्तत्र युध्यमानो मुहुर्मुहुः}
{नष्टाः सर्वे सुरास्तेन निर्जितो मधुसूदनः}% ॥७८॥

\twolineshloka
{निर्जितं तेन दैत्येन बाहुयुद्धमजायत}
{बाहुयुद्धं कृतं तेन दिव्यं वर्षसहस्रकम्}% ॥७९॥

\twolineshloka
{विष्णुश्चिन्तां प्रपन्नश्च नष्टाः सर्वाश्च देवताः}
{विष्णुश्च निर्जितस्तेन गतो बदरिकाश्रमम्}% ॥८०॥

\twolineshloka
{गुहा सिंहावती नाम तत्र सुप्तो जनार्दनः}
{योजनद्वादशवती एकद्वारा च पाण्डव}% ॥८१॥

\twolineshloka
{तस्यां विष्टः प्रसुप्तश्च दानवो हन्तुमुद्यतः}
{महायुद्धेन तेनैव श्रान्तोऽसौ योगमायया}% ॥८२॥

\twolineshloka
{दानवः पृष्ठतो लग्नो प्रविष्टः स तदा गुहाम्}
{प्रसुप्तं तत्र मां दृष्ट्वा दानवो हर्षमागतः}% ॥८३॥

\twolineshloka
{इत्थं मां निर्जितं मत्वा प्रविष्टं शङ्कया हरिम्}
{निःसन्देहं हनिष्यामि दानवानां भयङ्करम्}% ॥८४॥

\twolineshloka
{निर्गता कन्यका तत्र विष्णुदेहाद्युधिष्ठिर}
{रूपवती सुसौभाग्या दिव्यप्रहरणायुधा}% ॥८५॥

\twolineshloka
{तस्य तेजोंऽशसम्भूता महाबलपराक्रमा}
{दृष्टा सा दानवेन्द्रेण मुरनाम्ना धनञ्जय}% ॥८६॥

\twolineshloka
{युद्धं समाहितं तेन स्त्रिया चैव प्रयाचितम्}
{कन्यका युध्यते तत्र सर्वयुद्धविशारदा}% ॥८७॥

\twolineshloka
{हुङ्कारैर्भस्मसाज्जातो मुर नामा महासुरः}
{निहते दानवे तस्मिंस्तत्र देवस्त्वबुध्यत}% ॥८८॥

\twolineshloka
{पतितं दानवं दृष्ट्वा ततो विस्मयमागतः}
{केनायं च हतो रौद्रो ह्यत्युग्रो मम शात्रवः}% ॥८९॥

\onelineshloka*
{अत्युग्रं च कृतं कर्म मम कारुण्यतः कृतम्}

\uvacha{कन्योवाच}
\onelineshloka
{तेन देवाश्च गन्धर्वा सयक्षोरगराक्षसाः}% ॥९०॥

\twolineshloka
{इन्द्राद्याः सकला जित्वा स्वर्गाच्चैव निराकृताः}
{हरिः सुप्तो मया दृष्टो मुरः पृष्ठे समागतः}% ॥९१॥

\twolineshloka
{संहरिष्यति त्रैलोक्यं सुप्ते चैव जर्नादने}
{तस्यास्तद्वचनं श्रुत्वा विष्णुर्वचनमब्रवीत्}% ॥९२॥
\onelineshloka*
{अहं च निर्जितो येन कथं सोऽपि त्वया जितः}

\uvacha{एकादश्युवाच}
\onelineshloka
{त्वत्प्रसादाच्च भो स्वामिन् महादैत्यो मया हतः}% ॥९३॥

\uvacha{श्रीभगवानुवाच}

\threelineshloka
{आनन्दं त्रिषु लोकेषु मुनयो देवता गताः}
{ब्रूहि त्वं च मया भद्रे यत्ते मनसि रोचते}
{ददामि च न सन्देहो यत्सुरैरपि दुर्लभम्}% ॥९४॥

\uvacha{एकादश्युवाच}

\twolineshloka
{यदि तुष्टोऽसि मे देव सत्यमुक्तं जनार्दन}
{वरमेकं तु वाञ्छामि हृदये च जगत्पते}% ॥९५॥

\twolineshloka
{प्रार्थयामि च देवेश ईप्सितं च मया प्रभो}
{यदि सत्यं जगन्नाथ तिस्रो वाचो ददासि मे}% ॥९६॥

\uvacha{श्रीभगवानुवाच}

\twolineshloka
{सत्यं सत्यं मया प्रोक्तं अवश्यं तव सुव्रते}
{तिस्रो वाचो मया दत्ता न चावाक्यं भवेदिह}% ॥९७॥

\uvacha{एकादश्युवाच}

\twolineshloka
{त्रिभुवनेषु च देवेश चतुर्युगेषु साम्प्रतम्}
{त्रिषु लोकेषु सर्वत्र तादृशं कुरु मे प्रभो}% ॥९८॥

\twolineshloka
{सर्वतीर्थप्रधानं हि सर्वविघ्नविनाशिनी}
{सर्वसिद्धिकरी देवी त्वत्प्रसादाद् भवाम्यहम्}% ॥९९॥

\twolineshloka
{मामुपोष्यन्ति ये भक्त्या तव भक्त्या जनार्दन}
{सर्वसिद्धिर्भवेत्तेषां यदि तुष्टोऽसि मे प्रभो}% ॥१००॥

\twolineshloka
{उपवासं च नक्तं च एकभक्तं करोति च}
{तस्य वित्तं च धर्मं च मोक्षं वै देहि माधव}% ॥१०१॥

\uvacha{विष्णुरुवाच}

\twolineshloka
{यत्त्वं वदसि कल्याणि तत् सर्वं च भविष्यति}
{सर्वान् मनोरथान्भद्रे दास्यसि त्वं च नान्यथा}% ॥१०२॥

\twolineshloka
{मम भक्ताश्च ये लोके ये च भक्तास्तु कार्त्तिके}
{चतुर्युगेषु विख्यातास्त्रिषु लोकेषु वै प्रभो}% ॥१०३॥

\twolineshloka
{त्वां च शक्तिमहं मन्ये एकादशीव्रतस्थिताः}
{मम पूजां करिष्यन्ति मोक्षगास्ते न संशयः}% ॥१०४॥

\threelineshloka
{तृतीया चाष्टमी चैव नवमी च चतुर्दशी}
{एकादशी विशेषेण तिथिरेषा हरिप्रिया}
{सर्वतीर्थाधिकं पुण्यं सत्यं सत्यं न संशयः}% ॥१०५॥

\twolineshloka
{इदं दत्त्वा वरं तस्यै तिस्रो वाचो न संशयः}
{हृष्टा पुष्टा च सञ्जाता एकादशी महाव्रता}% ॥१०६॥

\twolineshloka
{शत्रुं हंसि परां तस्य ददासि परमां गतिम्}
{त्वं हंसि सर्वविघ्नानि सर्वसिद्धिवरप्रदा}% ॥१०७॥

\twolineshloka
{उभयोः पक्षयोः पार्थ तुल्या एकादशी शुभा}
{न शुक्ला नैव कृष्णा च विभेदं नैव कारयेत्}% ॥१०८॥

\twolineshloka
{विभेदो नैव कर्तव्यः समस्तव्रतकारिभिः}
{दिवा वा यदि वा रात्रौ शृणोति भक्तितत्परः}% ॥१०९॥

\twolineshloka
{तिथिरेका भवेत् सर्वा पक्षयोरुभयोरपि}
{उभयैकादशी स्वल्पा ह्यन्ते चैव त्रयोदशी}% ॥११०॥

\twolineshloka
{मध्ये तु द्वादशी पूर्णा त्रिस्पृशा सा हरिप्रिया}
{एकामुपोषयेत्तां वै सहस्रैकादशीफलम्}% ॥१११॥

\twolineshloka
{सहस्रगुणितं ह्येवं द्वादश्यां पारणे कृते}
{अष्टम्येकादशी षष्ठी तृतीया च चतुर्दशी}% ॥११२॥

\twolineshloka
{पूर्वविद्धा न कर्तव्या परविद्धामुपोषयेत्}
{एकादशी ह्यहोरात्रं प्रभाते घटिका भवेत्}% ॥११३॥

\twolineshloka
{सा तिथिः परिहर्तव्या उपोष्या द्वादशीयुता}
{एवंविधा मया प्रोक्ता पक्षयोरुभयोरपि}% ॥११४॥

\twolineshloka
{एकादश्यां प्रकुर्वीत ह्युपवासं न संशयः}
{ते यान्ति वैष्णवं स्थानं यत्राऽऽस्ते गरुडध्वजः}% ॥११५॥

\twolineshloka
{धन्यास्ते मानवा लोके विष्णुभक्तिपरायणाः}
{एकादश्याश्च माहात्म्यं सर्वकालेषु यः पठेत्}% ॥११६॥

\twolineshloka
{गोसहस्रफलं सोऽपि पुण्यं प्राप्नोति मानवः}
{दिवा वा यदि वा रात्रौ ये वै शृण्वन्ति भक्तितः}% ॥११७॥

\threelineshloka
{ब्रह्महत्यादिपापेभ्यो मुच्यन्ते नात्र संशयः}
{विष्णुधर्मसमं नास्ति गीतार्थं च नृपोत्तम}
{एकादशीसमं नास्ति व्रतं पापप्रणाशनम्}% ॥११८॥

॥इति श्रीपाद्मे महापुराणे पञ्चपञ्चाशत्साहस्र्यां संहितायामुत्तरखण्डे उमापतिनारदसंवादान्तर्गतकृष्णयुधिष्ठिरसंवादे मार्गशीर्ष-कृष्ण-उत्पन्ना-एकादशी-कृत-मुरवधो नाम चत्वारिंशोऽध्यायः॥४०॥

\genMangalaShloka

\hyperref[sec:ekadashi_mahatmyam_padma_puranam]{\closesub}
\clearpage

